\section{Introduction}
\label{sec:intro}

Visual odometry and \ac{lidar}-based \ac{slam} are well-established, but both
degrade in dusty, foggy, featureless, or textureless environments.
Millimeter-wave (mmWave) \ac{fmcw} radar is inherently robust to these
conditions: it penetrates smoke and dust, operates in the dark, and
measures Doppler velocity directly without feature matching.  These
properties make it attractive for aggressive flight in unstructured
environments.

Combining a mmWave radar with an \ac{imu} (\ac{rio}) for ego-motion estimation
has been explored with discrete-time \ac{ekf} and graph-optimization
formulations~\cite{doer2020ekf,kramer2020rve}.  However, the
asynchronous nature of radar-\ac{imu} data (radar at $\sim$11\,Hz, \ac{imu} at
$\sim$1\,kHz) motivates a continuous-time representation where the
trajectory is a smooth function evaluated at any sensor timestamp
without interpolation approximation.

\textbf{Research question.}
This work asks, end to end: \emph{can a single-chip mmWave radar,
fused with an \ac{imu}, provide usable 6-DoF state estimation --- and if
so, what does it take?}  The starting point is what the sensor can and
cannot do: per-point Doppler yields direct ego-velocity, but at
10--60 points per frame there are no bearing-stable features for
position fixing, and yaw is a gauge freedom of Doppler+\ac{imu}.  The
system is therefore a \emph{velocity and orientation} source whose
position estimate drifts with distance --- by construction, not by
defect.  What it takes to turn that capability into accurate
estimates, and where it ultimately fails, is the subject of the rest
of this paper.

We contribute:
\begin{itemize}
  \item \emph{Single-chip radar under aggressive and tumbling flight},
        characterizing how the radar measurement model breaks in this
        regime: rate-dependent intra-frame smearing, and the finding
        that the apparent elevation bias under tumbling is a
        flight-regime error proxy, not a physical antenna constant.
        Backflips at \SI{10}{\radian\per\second} on a single-chip
        sensor are, to our knowledge, unexplored.
  \item A \emph{consistency-analyzed} fixed-lag smoother: a factor-set
        selection rule exploiting B-spline local support to make the
        Schur marginal exactly closed --- distinct from and
        complementary to \ac{fej} --- eliminating prior-weight tuning, at
        0.35--0.70\,s per window on a laptop (real time for aggressive
        racing at a 0.4\,s stride).
  \item One measurement-weighting configuration from hover to
        \SI{10}{\radian\per\second}, with the velocity--orientation
        trade as a measured Pareto curve and accuracy validated on
        \emph{held-out flights} never used for tuning.
  \item Honest negative results: a diagnostic symptom pattern for
        marginalization inconsistency; a representation-level study
        that \emph{disproves} the adaptive-knot-density hypothesis for
        this regime; ground-truth limits during aggressive maneuvers;
        and a \ac{nees} check confirming the reported orientation covariance
        is consistent on racing flights.
\end{itemize}

\textbf{Data and code.}
The complete system --- physically correct Doppler model with lever
arm, sensor-only initialization, full hyperparameter disclosure, and
deterministic re-runs --- is released for independent reproduction.
The datasets (Vicon \ac{mocap} lab, TI IWR6843AOPEVM, Pixhawk) and the
C++ Ceres solver with Python bindings are available at
\url{https://github.com/Dr1zz3l/radar-iwr6843-driver}.
